\chapter{Introduction}

Logic minimization is an essential step in low-level synthesis. The most important parameters for evaluating a logic circuit are its physical area and processing speed. A fundamental unit, such as the universal \ac{NAND} gate, is built using a complementary arrangement of multiple transistors (typically two \ac{NMOS} and two \ac{PMOS}) and can be used to construct any logic circuit. Each logic gate requires a specific amount of time to process a signal and occupies physical space. Consequently, a logic circuit is implemented using various interconnected logic blocks, all of which contribute to the total size and execution time. A logic function can be described as a \ac{SOP}, mapping directly to a \ac{PLA} structure comprising an AND plane and an OR plane. Minimizing the number of literals and product terms reduces the required logic gate count, ultimately yielding a faster and smaller circuit.

Any form of logic minimization is \ac{NP}-hard, meaning no known polynomial-time algorithm exists to find the optimal solution. However, with heuristics, that is, obtaining a solution that works well enough, the problem becomes tractable. That is where the Espresso algorithm comes in.

Espresso was developed by Brayton, Hachtel, McMullen, Rudell and Sangiovanni-Vincentelli at UC Berkeley (1984)\cite{demicheli1994}, and is a standard tool used for logic minimization.
The Espresso algorithm is a heuristic-based logic minimization tool that iteratively optimizes a two-level \ac{SOP} representation. Its objective is to group the fixed mathematical minterms of a function into the most efficient cover possible. It achieves this by simultaneously minimizing the total number of product terms (cubes) to reduce the required number of logic gates, and minimizing the number of literals per term to reduce the transistor count per gate, without altering the underlying boolean logic. 

Internally, Espresso operates on three disjoint sets - the ON set, OFF set, and don't-care set (DC set). Espresso iterates through three phases: expand, which enlarges each cube as far as possible without covering the OFF-set; irredundant, which removes redundant cubes from the cover; and reduce, which shrinks cubes to enable better expansions in the next iteration. This loop repeats until no further reduction in the cost function is achievable.

The Espresso implementation presented in the report includes complement set calculation, a cost function defined as a cover size, and three phases of core Espresso algorithm - Expand, irredundant, and Reduce, which iteratively minimizes cost function. The Espresso implmentation is performed in Java, with logic function equivalency checked using ABC \cite{abc}. Each part's implementation is explained in section \textit{Implementation}.\ainote{Claude Sonnet 4.6}{Improve grammar and LaTeX formatting of the report.}